\section{Related work}
\label{sec:related}
\subsection{Origami benchmarks and generation}
\label{sec:related-benchmarks}
GamiBench evaluates visual questions about folded configurations and crease patterns
\citep{gamibench}. OrigamiSpace covers prediction, spatial relationships and crease-pattern
code generation \citep{origamispace}; CP2Seq instead supplies the pattern and asks for a
verified folding sequence. OrigamiBench synthesises flat-foldable patterns and scores
query efficiency, image overlap and CLIP similarity, without verifying final layer order
\citep{origamibench}. Corigami uses geometric and semantic evaluation for generation
\citep{corigami}, and Learn2Fold plans against a learned world model \citep{learn2fold}.
FoldingAgent reconstructs sequences from instructional video \citep{foldingagent}, using
the human-authored PurelandFold corpus \citep{purelandfold}; CP2Seq neither trains nor
evaluates on it. Our targets are generated by recorded legal sequences, and final stacks
are checked layer by layer and by original-sheet region.

\subsection{Computational origami and sequence planning}
\label{sec:related-theory}
Simple folds operate along one line at a time; single-, some- and all-layers models
restrict which layers move \citep{arkin-map}. Foldability complexity depends on paper
shape, crease directions and permitted operations \citep{akitaya-hard,akitaya-infinite,mixed-orthogonal}.
Layer ordering is central to flat-folded states
\citep{bern-hayes,akitaya-flatfolder,layer-algebra,flat-folding-graphs}.
Continuous reachability results \citep{continuous,reaching-folded} do not establish
reachability by a restricted sequence of simple folds. Sequence recovery from crease
patterns is an established computational problem \citep{akitaya-cp2seq}; CP2Seq applies
it to language-model planning with executable actions and matched search baselines.
Because each target has a legal reference sequence, benchmark reachability is promised,
and general foldability hardness does not establish hardness of this generated distribution
(Appendix~\ref{app:folding-complexity}).
